\documentclass[12pt,a4paper]{article}

\usepackage[utf8]{inputenc}
\usepackage[T1]{fontenc}
\usepackage{lmodern}
\usepackage[english]{babel}
\usepackage{csquotes}
\usepackage{amsmath,amssymb}
\usepackage{booktabs}
\usepackage{tabularx}
\usepackage{adjustbox}
\usepackage{graphicx}
\usepackage{xcolor}
\usepackage[hypertexnames=false]{hyperref}
\usepackage{cleveref}
\usepackage{float}

\usepackage[style=apa,backend=biber,natbib=true]{biblatex}
\addbibresource{references.bib}
\DeclareLanguageMapping{english}{english-apa}

\newcommand{\xscore}{\textit{xScore}}
\newcommand{\xvoa}{\textit{xVOA}}
\newcommand{\gds}{\textit{GDS}}

\emergencystretch=2em
\widowpenalty=10000
\clubpenalty=10000

\title{Online Supplement: Game Deserved Score --\\Decomposing NFL Team Quality and Testing\\the Defensive Dominance Hypothesis}
\author{Lasse Mecha}
\date{2026}

\begin{document}
\maketitle

\noindent This supplement reports validation and robustness evidence referenced in, but not
reproduced in, the main paper. All numbers use the same 256-team-season sample (2018--2025)
and the same fixed \xscore{}/\gds{} pipeline as the main paper; none of the material here
changes the paper's conclusions, which is why it appears separately rather than in the body.
Section numbers below are independent of the main paper's numbering and are cross-referenced
from it by section title, not by number.

\section*{S1. Luck Analysis}

\gds{}-implied wins are the fitted values from the season-level regression of win percentage
on \gds{}/game, estimated on the full 256-team-season sample and applied to the most recently
completed season (2025). The gap between actual and implied wins operationalizes luck.

\begin{table}[htbp]
  \centering
  \caption{Luckiest and unluckiest teams in the 2025 NFL regular season by divergence between
  actual wins and \gds{}-implied wins, using the full-sample (2018--2025) regression fit.}
  \begin{tabular}{lrrr}
    \toprule
    \textbf{Team} & \textbf{Actual Wins} & \textbf{GDS-Implied} & \textbf{Divergence} \\
    \midrule
    \multicolumn{4}{l}{\textit{Luckiest}} \\[2pt]
    CHI & 11 & 8.5  & $+2.5$ \\
    DEN & 14 & 11.7 & $+2.3$ \\
    PIT & 10 & 7.8  & $+2.2$ \\[6pt]
    \multicolumn{4}{l}{\textit{Unluckiest}} \\[2pt]
    NYG &  4 & 6.8  & $-2.8$ \\
    IND &  8 & 11.2 & $-3.2$ \\
    KC  &  6 & 9.7  & $-3.7$ \\
    \bottomrule
  \end{tabular}
\end{table}

Kansas City is the most striking case: moderate \gds{} implying roughly 9.7 wins, yet only 6
actual wins -- the largest negative divergence in the sample. Chicago, Denver, and Pittsburgh
each outperformed their process-quality expectation by 2+ wins. These divergences reinforce the
main paper's core motivation for using \gds{} rather than win-loss record alone: win totals are
an imperfect proxy for underlying quality, and process-based measurement captures information
outcome-based records do not.

\section*{S2. Baseline Comparison: Point Differential}

Raw point differential per game is trivially available and requires no model. Across the full
256-team-season sample, point differential achieves $r = 0.900$ with win percentage
($R^2 = 81.0\%$), compared to \gds{}'s $r = 0.862$ ($R^2 = 74.3\%$) -- a 6.7-percentage-point
gap. This is expected and does not undermine \gds{}'s utility: point differential contains the
very outcome information \gds{} deliberately excludes (it is partially tautological, since wins
are defined by outscoring the opponent), whereas \gds{} measures process quality net of that
outcome variance. The value of \gds{} over point differential is not raw predictive power but
three properties point differential lacks: decomposability into offense/defense/special-teams
components, process measurement that treats garbage-time execution differently from
competitive-situation execution (via score-differential conditioning), and separation of skill
from the outcome variance that S1's luck analysis exploits directly.

\section*{S3. Per-Season Correlation Stability}

\begin{table}[htbp]
  \centering
  \caption{Per-season Pearson $r$ between \gds{}/game and win percentage ($n = 32$ teams per
  season, all eight seasons in the sample).}
  \begin{tabular}{lcc}
    \toprule
    \textbf{Season} & \textbf{Pearson $r$} & \textbf{$R^2$} \\
    \midrule
    2018 & 0.871 & 75.9\% \\
    2019 & 0.818 & 67.0\% \\
    2020 & 0.881 & 77.7\% \\
    2021 & 0.877 & 76.9\% \\
    2022 & 0.886 & 78.6\% \\
    2023 & 0.847 & 71.7\% \\
    2024 & 0.875 & 76.5\% \\
    2025 & 0.888 & 78.8\% \\
    \midrule
    \textit{Pooled} & 0.862 & 74.3\% \\
    \bottomrule
  \end{tabular}
\end{table}

Per-season correlations range from $r = 0.818$ (2019) to $r = 0.888$ (2025 -- the newest,
out-of-sample season, at the high end of the range rather than an outlier), with no season
below $r = 0.82$. This confirms the pooled correlation is not driven by any single season and
that the relationship between \gds{} and winning is stable across the entire study window,
including the season added after the underlying model was built.

\section*{S4. Year-Over-Year Rank Stability}

Across the seven consecutive-season pairs available in an eight-season sample (2018--2019
through 2024--2025), the Spearman rank correlation between consecutive-season \gds{}/game
rankings averages $\rho = 0.374$ ($n = 32$ per pair; range $0.183$ to $0.617$):
2018--19: $\rho=0.617$; 2019--20: $\rho=0.489$; 2020--21: $\rho=0.479$; 2021--22: $\rho=0.207$;
2022--23: $\rho=0.352$; 2023--24: $\rho=0.183$; 2024--25: $\rho=0.289$. This moderate
persistence is consistent with the NFL's competitive-balance mechanisms (draft order, salary
cap, free agency) promoting substantial year-over-year turnover while not eliminating quality
persistence entirely.

\section*{S5. Split-Half Reliability}

Correlating \gds{}/game computed from odd-numbered weeks against even-numbered weeks tests
whether roughly half a season's data produces consistent team rankings. The Spearman rank
correlation, computed separately for each season and averaged, is $\rho = 0.477$ (per-season:
2018: 0.474; 2019: 0.356; 2020: 0.632; 2021: 0.580; 2022: 0.302; 2023: 0.413; 2024: 0.560; 2025:
0.500). The Spearman-Brown prophecy formula implies a full-season reliability of
$2 \times 0.477 / (1 + 0.477) = 0.646$, indicating adequate measurement reliability for
cross-team comparison at full-season aggregation, with meaningful differentiation available by
roughly mid-season.

\section*{S6. Full Robustness Detail}

Summary versions of all six checks appear in the main paper's robustness table; full detail is
reported here.

\paragraph{Opponent-strength control.} Re-estimating the OLS model (playoff wins on offensive
and defensive \xvoa{}/game) with a leave-one-out opponent-quality covariate added:
$\hat\beta_{\text{off}} = 0.099$ ($p<0.0001$), $\hat\beta_{\text{def}} = 0.071$ ($p<0.0001$),
$\hat\beta_{\text{opp}} = -0.013$ ($p = 0.830$), $R^2 = 0.267$ (unchanged from the uncontrolled
model). The opponent-strength coefficient is statistically indistinguishable from zero and
neither offensive nor defensive coefficients move, providing direct evidence against the
schedule-strength confound.

\paragraph{Field-position mediation.} The correlation between defensive \xvoa{}/game and a
team's average starting field-position advantage is $r = -0.151$ ($p = 0.0157$) -- a real but
weak mediating pathway. The total correlation between offensive \xvoa{}/game and playoff wins
is $r = 0.439$; the partial correlation controlling for starting field position is
$r_{\text{partial}} = 0.404$. Mediation percentage: $(0.439 - 0.404)/0.439 = 8.0\%$. Just over
90\% of the offensive-\xvoa{}-to-wins association is independent of any field-position benefit
that better defenses might confer on their own offense.

\begin{table}[htbp]
  \centering
  \caption{Threshold sensitivity: mean playoff wins for defense-dominant team-seasons across
  the full range of archetype-threshold definitions tested.}
  \begin{tabular}{lcc}
    \toprule
    \textbf{Threshold} & \textbf{$n$ (Def-Dominant)} & \textbf{Mean PW} \\
    \midrule
    $\pm0.20$ & 42 & 0.000 \\
    $\pm0.25$ & 38 & 0.000 \\
    $\pm0.30$ & 32 & 0.000 \\
    $\pm0.35$ & 28 & 0.000 \\
    $\pm0.40$ & 26 & 0.000 \\
    \bottomrule
  \end{tabular}
\end{table}

\begin{table}[htbp]
  \centering
  \caption{Time-window sensitivity of the offensive-to-defensive $R^2$ ratio.}
  \begin{tabular}{lcc}
    \toprule
    \textbf{Window} & \textbf{$R^2$ Ratio (Off:Def)} & \textbf{$n$ (team-seasons)} \\
    \midrule
    2018--2025 (full) & 2.6:1 & 256 \\
    2018--2023        & 3.2:1 & 192 \\
    2020--2025        & 2.8:1 & 192 \\
    2021--2025        & 2.7:1 & 160 \\
    2018--2021        & 2.5:1 & 128 \\
    \bottomrule
  \end{tabular}
\end{table}

\paragraph{Single-team exclusion.} Excluding Kansas City 2023 (lowest offense share among
champions before this paper's extension): ratio $= 2.7$:$1$, Spearman $\rho = 0.221$
($p = 0.022$). Excluding Seattle 2025 (the newest champion, to confirm the newest season alone
is not driving the result): ratio $= 2.7$:$1$, Spearman $\rho = 0.215$ ($p = 0.026$). Both
exclusions leave the full-sample estimates ($2.6$:$1$, $\rho = 0.208$) essentially unchanged.

\begin{table}[htbp]
  \centering
  \caption{Era comparison. Spearman: offense share vs.\ playoff wins (playoff teams only).
  Pearson: \xvoa{} components vs.\ win percentage (all team-seasons).}
  \begin{tabular}{lrrr}
    \toprule
    \textbf{Measure} & \textbf{2018--2020} & \textbf{2021--2025} & \textbf{$\Delta$} \\
    \midrule
    Spearman $\rho$ (share vs.\ PW)      & $0.305$ ($p=0.063$) & $0.163$ ($p=0.177$) & $-0.142$ \\
    Sample ($n$, playoff teams)          & 38                  & 70                   & \\[4pt]
    Pearson $r$ (off \xvoa{} vs.\ Win\%) & $0.622$             & $0.701$              & $+0.079$ \\
    Pearson $r$ (def \xvoa{} vs.\ Win\%) & $0.397$             & $0.427$              & $+0.030$ \\
    \bottomrule
  \end{tabular}
\end{table}

Neither era's Spearman coefficient reaches significance individually, and -- as the main paper
notes explicitly in its Structural Explanations discussion -- the defensive Pearson coefficient
strengthens slightly in the later era rather than weakening, which is different from what the
224-team-season-only version of this analysis showed. We report the updated numbers as-is
rather than retaining the earlier, cleaner-sounding pattern.

\section*{Data and Code Availability}

All numbers in this supplement are reproducible from
\texttt{papers/paper2-merged/compute\_placeholders.py}, run against
\texttt{output/archetype\_v2\_data.csv} and \texttt{output/multinomial\_game\_gds.csv}. See the
main paper's Data and Code Availability statement for the full repository link.

\printbibliography

\end{document}
